\chapter{Evaluation}\label{ch:eval}

\section{Evaluation goals and methodology}

This chapter evaluates the core goals of Veritas put forward in the proposal and refined in the requirements analysis (Section~\ref{sec:requirements-anal}).
A central claim of Veritas' architecture is that low-level artefacts can be independently verified while the compiler can be untrusted. Veritas also aimed to support a toolchain for a \emph{systems language}. The evaluation therefore has to assess both the verification story and the ordinary compiler story:

\begin{itemize}
  \item Does DTAL realise the intended trust architecture?
  \item Is the implementation a viable systems-language toolchain?
  \item Is the source language expressive enough for the targeted verification tasks?
  \item Is independent verification practical enough to use?
  \item How much confidence should we have in the verifier?
  \item How much does the architecture reduce the TCB?
\end{itemize}

The evidence used for these questions is deliberately mixed: comparison with Xi and Harper's DTAL, feature-suite coverage, invalid-program rejection, DTAL tampering tests, native code generation, PolyBench-style kernels, wall-clock measurements, solver breakdowns, and a trusted-computing-base analysis for the Linux and bare-metal targets.

The evaluation uses fixed suites rather than ad hoc examples: 28 valid feature programs, 35 invalid source programs, 14 syntax-preserving DTAL tampering cases, and six cross-system comparison tasks. Each suite has a different evidential role. The positive feature suite checks that representative supported programs make it through both the frontend and the independent verifier. The negative suite checks that unsafe source programs are rejected rather than accidentally accepted. The DTAL tampering corpus starts from valid compiled DTAL and mutates the low-level program, directly testing the untrusted-compiler story. The cross-system suite gives a small normalised comparison against established verification tools.

Timing results use two measurement layers. Internal metrics from \texttt{veritas --bench} separate source-to-DTAL compilation, explicit DTAL verification, and SMT time; \texttt{hyperfine} wall-clock timings include process startup and native-output generation. Runtime correctness is checked using exact outputs or deterministic checksums, not successful termination alone. This distinction is important because a verifier can appear cheap or expensive depending on whether one measures only the checking phase, the full command invocation, or the generated program's runtime behaviour. The chapter therefore treats these as separate claims rather than folding them into one performance number.

\section{DTAL design and implementation}\label{sec:dtal-eval}

In the Veritas architecture, the target DTAL must be expressive enough to enable verification of the required safety properties. The evaluation considers to what extent the DTAL implementation supports the intended trust architecture.

\subsection{Relationship to Xi and Harper's DTAL}

Veritas is inspired directly by Xi and Harper's DTAL and preserves many of its central features and ideas most relevant for independent low-level verification:

\begin{center}
\small
\begin{tabularx}{0.96\textwidth}{>{\raggedright\arraybackslash}p{0.32\textwidth} >{\raggedright\arraybackslash}X}
\toprule
\textbf{DTAL mechanism} & \textbf{Safety role} \\
\midrule
Dependent block or label state & Supports control-flow safety by ensuring that jumps and joins are only valid when the incoming machine state matches the target's expected typing context. \\
Singleton integer typing & Supports arithmetic and indexing safety, since later instructions can rely on symbolic facts when proving side conditions. \\
Size-indexed arrays & Supports memory safety by making out-of-bounds reads and writes statically checkable. \\
Typed stack discipline & Supports structural safety by preventing ill-typed stack use across calls, returns, and local temporaries. \\
Proof-relevant typed memory access & Supports static memory-safety guarantees such as array-bound-check elimination. \\
\bottomrule
\end{tabularx}
\end{center}

Veritas narrows Xi and Harper's design in some directions. It does not expose their full tuple and allocation primitives, broad existential dependent types, or sum and tag reasoning. Instead, it focuses on the low-level facts needed for fixed-size arrays, arithmetic side conditions, contracts, and borrow-aware compilation. This sacrifices theoretical generality in exchange for automation and a verification model that fits the rest of the toolchain.

The most important continuity with Xi and Harper's DTAL lies in the independent low-level verification model. The Veritas version of the copy-loop pattern from Section~\ref{sec:dtal-background} preserves the same basic idea of typed labels and register-indexed array access, although the surface syntax is adapted to its implementation:

\begin{quote}
  \begin{verbatim}
.entry {v0: int}
.assume (v0 >= 0 && v0 < 4)
add v9, v0, v8    : {_synth_0: int | _synth_0 == (i + 1) }
.assert true
load v10, [v1 + v9]    : int(7)
\end{verbatim}
\end{quote}

The low-level program records sufficient information for the verifier to justify array access from the surrounding type and constraint context rather than relying on a dynamic bounds check.

\subsection{Extensions and deviations}


Veritas preserves the same basic typed memory access idea, but also extends it with explicit refinement-carrying annotations that originate in source-level proof obligations, for example:
\begin{quote}
  \begin{verbatim}
div v2, v0, v1    : {_synth_0: int | _synth_0 == (a / b) }
load v10, [v1 + v9]    : int(7)
  \end{verbatim}
\end{quote}
These carried obligations are re-checked by the verifier against the low-level code. This provides a broader source-to-low-level continuity for the particular fragment Veritas supports, including not just bounds facts but also contract-style and borrow-related safety obligations. Preserving additional proof-relevant annotations does increase the surface on which an untrusted compiler can introduce inconsistencies, and therefore may increase the risk of false rejection, but it does not expand the acceptance surface or weaken verification guarantees.

Another extension is the introduction of explicit ownership and borrow-aware operations.
This extension pushes alias-control properties below the frontend where the relevant low-level operations are represented and re-checked in DTAL. This is a significant divergence from Xi and Harper's DTAL in order to serve a systems-language setting.

\begin{quote}
  \begin{verbatim}
alias_borrow v3, v1    : &[int(7); 1]
load v5, [v3 + v4]     : int
borrow_end v3          : &[int(7); 1]
  \end{verbatim}
\end{quote}

In this example, the creation, use, and termination of the borrowed view is represented explicitly in DTAL, which allows the verifier to check that the access occurs through a live borrow and that the borrow is ended consistently with the source-language ownership discipline.

Because verification happens after physical register allocation, Veritas defines a `physical DTAL' with explicit scratch-register use, spills, prologues/epilogues, and instruction expansions. Xi and Harper's DTAL is closest to Veritas' virtual DTAL: its semantics are given over an abstract machine, and register allocation and spilling are explicitly not modelled. The difference is visible in instructions such as:

\begin{quote}
  \begin{verbatim}
mov lr, r1    : int
cqo
idiv r7
mov r0, lr    : int
\end{verbatim}
\end{quote}

or in explicit stack and calling-convention machinery such as
\texttt{spill\_store}, \texttt{spill\_load}, \texttt{.prologue}, and
\texttt{.epilogue}.
The main tradeoff is low-level complexity. Physical DTAL exposes register assignments, spill slots, calling-convention details, and instruction expansion that are absent from the source language. This makes the representation less elegant than an idealised typed assembly and increases the burden on the verifier. However, it also makes register allocation and instruction expansion part of the untrusted compiler, which is more useful for the project's trust model.




\section{Veritas as a systems-language toolchain}\label{sec:veritas-sys-eval}

This section evaluates Veritas as a systems-language toolchain. The project applies the trust architecture in a systems-oriented setting, where memory safety has to coexist with imperative control flow, mutable state, references, predictable data representation, and native code generation. Veritas is far less mature than C or Rust in language breadth, tooling, and ecosystem, so the relevant question is narrower: within the intended fragment, does it behave like a plausible systems-language design that could be built on further?


\subsection{Systems-language feature comparison}

Compared with C and Rust, Veritas occupies an intermediate position. Like both languages, it supports first-order functions, mutable local state, fixed-size arrays, loop-based imperative control flow, and native-code compilation. Like Rust, it has an ownership- and borrow-aware safety story rather than relying purely on unrestricted memory manipulation. Unlike both, its design is centred around carrying verification-relevant information through the compilation pipeline. This makes Veritas less general-purpose than C or Rust, but gives it a stronger explicit safety and trust architecture within its supported fragment.


\begin{table}[H]
  \centering
  \small
  \begin{tabularx}{0.96\textwidth}{>{\raggedright\arraybackslash}X c c >{\raggedright\arraybackslash}p{0.22\textwidth}}
    \hline
    \textbf{Feature}                            & \textbf{C} & \textbf{Rust} & \textbf{Veritas}                      \\
    \hline
    Mutable local state                         & \checkmark & \checkmark    & \checkmark                            \\
    Fixed-size arrays                           & \checkmark & \checkmark    & \checkmark                            \\
    Structured loop control flow                & \checkmark & \checkmark    & \checkmark                            \\
    Native-code compilation                     & \checkmark & \checkmark    & \checkmark                            \\
    Built-in ownership discipline               &            & \checkmark    & \checkmark                            \\
    Refinement/contracts/invariants             &            &               & \checkmark                            \\
    Independent low-level verifier              &            &               & \checkmark                            \\
    User-defined structs / algebraic data types & \checkmark & \checkmark    & arrays / refs only                    \\
    General heap allocation model               & \checkmark & \checkmark    & restricted \texttt{alloca} only       \\
    General unrestricted pointers               & \checkmark & \checkmark    &                                       \\
    Basic I/O and runtime interaction           & \checkmark & \checkmark    & partial                               \\
    Mature module / package ecosystem           & \checkmark & \checkmark    &                                       \\
    FFI / systems-library integration story     & \checkmark & \checkmark    & minimal                               \\
    Heap-heavy general systems idioms           & \checkmark & \checkmark    & array / ref kernels only              \\
    Concurrency model in current fragment       & \checkmark & \checkmark    &                                       \\
    \hline
  \end{tabularx}

  \caption{High-level comparison of Veritas with C and Rust as systems-language
    toolchains.}
  \label{tab:systems-language-comparison}
\end{table}

The shared features are enough for the low-level kernels and array-oriented routines used throughout the evaluation. The missing features still limit its generality, but within a narrower fragment centred on structured imperative code with strong safety tracking, Veritas exhibits recognisably systems-language behaviour while offering a more explicit verification and trust model than conventional toolchains.

\subsection{Compiler functional correctness}

The architecture treats the compiler as untrusted, but it should still behave like a plausible toolchain. Since the compiler is not formally proved correct, this part of the evaluation is empirical and corpus-based: it asks whether generated binaries agree with the intended semantics on fixed suites.

\begin{center}
\small
\begin{tabularx}{0.96\textwidth}{>{\raggedright\arraybackslash}p{0.22\textwidth} >{\raggedright\arraybackslash}p{0.18\textwidth} >{\raggedright\arraybackslash}X}
\toprule
\textbf{Suite} & \textbf{Result} & \textbf{Role in the argument} \\
\midrule
\texttt{feature\_suite} & 28 / 28 pass & Local correctness corpus for arrays, mutable state, borrowing, contracts, invariants, quantifiers, safe division, and overflow-aware arithmetic. \\
PolyBench-style kernels & 7 / 7 pass & External check on recognised array- and loop-dominated workloads, validated by stable checksums against C baselines. \\
\bottomrule
\end{tabularx}
\end{center}

The feature suite checks the constructs and proof obligations Veritas was built to support. The PolyBench-style suite is more objective: PolyBench/C is widely used for compiler and program-transformation evaluation, and its kernels contain regular control flow, affine array indexing, and deterministic numeric behaviour.

\begin{table}[h]
  \centering
  \small
  \begin{tabularx}{0.96\textwidth}{>{\raggedright\arraybackslash}p{0.17\textwidth} >{\raggedright\arraybackslash}p{0.18\textwidth} >{\raggedright\arraybackslash}X}
    \hline
    \textbf{Kernel} & \textbf{\texttt{MEDIUM} params} & \textbf{Why it is included} \\
    \hline
    Floyd--Warshall & \texttt{N=500} & Dynamic programming over a weighted graph matrix; stresses nested loops, fixed-size 2D arrays, arithmetic updates, and indexed mutation. \\
    Jacobi-1D       & \texttt{N=400}, \texttt{T=100} & Iterative stencil update over repeated time steps; stresses affine indexing and arithmetic recurrences. \\
    Jacobi-2D       & \texttt{N=250}, \texttt{T=100} & Two-dimensional stencil computation; stresses neighbourhood indexing over fixed-size arrays. \\
    Seidel-2D       & \texttt{N=400}, \texttt{T=100} & In-place smoothing stencil; stresses array mutation and affine neighbour access. \\
    MVT             & \texttt{N=400} & Matrix-vector and transposed-matrix-vector multiplication; stresses regular matrix/vector indexing. \\
    ATAX            & \texttt{M=390}, \texttt{N=410} & Matrix-transpose-times-vector computation; stresses multiple loop phases and intermediate accumulation. \\
    GESUMMV         & \texttt{N=250} & Weighted sum of two matrix-vector products; stresses linear-algebra style arithmetic and nested loops. \\
    \hline
  \end{tabularx}
  \caption{PolyBench-style kernels used in the Veritas evaluation, with their
    fixed \texttt{MEDIUM}-class parameters and the main compiler behaviours
    they exercise.}
  \label{tab:polybench-kernels}
\end{table}

Table~\ref{tab:polybench-kernels} gives a spread across dynamic-programming, stencil, and matrix/vector-style computations. These are appropriate for Veritas because the language and verifier are strongest on fixed-size arrays, structured imperative loops, arithmetic side conditions, and deterministic computation over statically bounded data. All seven kernels pass their defined semantic checks.

While these results do not amount to full compiler correctness in the CompCert sense, nor does it imply that every future program in the Veritas fragment will automatically enjoy the same guarantee, it supports a strong suite-bounded claim: for the representative local feature corpus and the committed external kernel corpus, the delivered compiler both preserves verifier-relevant structure and produces binaries whose observed behaviour matches the intended semantics.

\subsection{Compiler and runtime practicality}

Practicality is evaluated at two levels: compilation cost and runtime behaviour of the generated binaries. The curated \texttt{feature\_suite} provides a useful summary of compiler throughput.

\begin{figure}[H]
  \centering
  \includegraphics[width=0.8\textwidth]{figures/chapter4_feature_suite_compile_cost.png}
  \caption{Compile-cost structure across the curated feature suite.}
  \label{fig:feature-suite-compile-cost}
\end{figure}


Figure~\ref{fig:feature-suite-compile-cost} makes the dominant cost pattern clear: binary-search, array-proof, and quantifier-heavy programs are the most expensive to compile, and a significant share of that cost is solver time. Simpler arithmetic- and borrow-oriented cases remain much cheaper.

The external PolyBench-style kernels provide a second view on larger loop- and array-heavy programs. Here compile time is dominated more strongly by frontend SMT work, because the kernels repeatedly exercise affine array indexing, nested loops, and arithmetic side conditions. Figure~\ref{fig:polybench-compile-cost} shows that the fastest kernel still compiles in under half a second, while the most expensive kernels take tens of seconds and are again dominated by solver-heavy proof structure rather than by raw source size alone.

\begin{figure}[H]
  \centering
  \includegraphics[width=0.92\textwidth]{figures/chapter4_polybench_compile_cost.png}
  \caption{Compile-cost structure across the translated PolyBench-style kernel suite.}
  \label{fig:polybench-compile-cost}
\end{figure}


Runtime practicality is evaluated using the same PolyBench-style kernel set to ask whether compiled verified kernels remain plausible executable programs. All included kernels pass semantic validation, and runtime performance relative to GCC varies substantially by workload.

\begin{figure}[H]
  \centering
  \includegraphics[width=0.92\textwidth]{figures/chapter4_polybench_runtime_absolute.png}
  \caption{Absolute runtime for the committed \texttt{MEDIUM}-class PolyBench-style suite. The figure compares Veritas against GCC \texttt{-O3} and checked C \texttt{-O3}; all plotted kernels passed checksum-based semantic validation before timing results were accepted.}
  \label{fig:polybench-runtime-absolute}
\end{figure}

The results in Figure~\ref{fig:polybench-runtime-absolute} show a mixed pattern rather than a uniform slowdown. On the \texttt{MEDIUM}-class PolyBench suite, MVT and GESUMMV are relatively close to the GCC \texttt{-O3} baseline, at roughly 1.55$\times$ and 1.10$\times$ respectively, while Floyd--Warshall, Jacobi-2D, and Seidel-2D are slower by roughly 3.3--5.9$\times$. This shows that runtime performance depends strongly on workload shape rather than following a single fixed overhead pattern.

To test whether the fixed-suite results were stable beyond one benchmark size, a follow-up size sweep was run on the more stable kernels. Each kernel was timed for Veritas, GCC \texttt{-O3}, and a checked C \texttt{-O3} baseline, where `checked C' denotes the manually instrumented variant with explicit dynamic array-bounds checks.

\begin{figure}[H]
  \centering
  \includegraphics[width=0.96\textwidth]{figures/chapter4_polybench_scaling_absolute.png}
  \caption{Runtime scaling for the chapter-facing PolyBench follow-up sweep. These size sweeps are anchored around the \texttt{MEDIUM}-class kernel sizes and compare Veritas with GCC \texttt{-O3} and checked C \texttt{-O3}.}
  \label{fig:polybench-scaling-absolute}
\end{figure}

Figure~\ref{fig:polybench-scaling-absolute} supports the same interpretation as the fixed-size suite. Stencil-style and dynamic-programming kernels remain slower than both C baselines, while the lighter matrix/vector kernels stay much closer. Veritas therefore does not incur one uniform slowdown on array-heavy programs; the relative runtime cost depends strongly on the computational pattern being compiled. The generated code is executable, semantically validated, and often within a small constant factor of compiled C on the tested kernels, while also participating in a verification pipeline that C itself does not provide.

Binary size is a supporting practicality metric. The larger loop-heavy kernels compile to small freestanding Veritas binaries, typically around 1.4--2.4\,KB. Veritas code is larger than GCC's \texttt{.text} section across the suite, but not by an order of magnitude, which is consistent with a research backend that prioritises verifier-visible structure over code density.

\section{Source language expressiveness for verification}\label{sec:source-lang-eval}

This section focuses on the source language as a verification-oriented language. It evaluates what kinds of properties it can express in practice, how that compares with related verification languages, and why the source-level expressiveness story is different in Veritas because the resulting proof-relevant information is preserved into DTAL rather than erased after frontend checking.

\subsection{Properties expressible in Veritas}

Within the supported fragment, Veritas can express contracts, loop invariants, bounded quantified properties, array-bounds and arithmetic safety conditions, sortedness-style properties, and borrowing constraints. It is not yet a mature verification ecosystem, but it is expressive enough for the dissertation's array-heavy programs.

\begin{center}
\small
\begin{tabularx}{0.96\textwidth}{>{\raggedright\arraybackslash}p{0.30\textwidth} >{\raggedright\arraybackslash}X}
\toprule
\textbf{Expressive strength} & \textbf{Current limit} \\
\midrule
Sorted arrays, bounded indices, loop invariants, arithmetic side conditions, postconditions, and ownership constraints. & No rich user-defined algebraic datatypes, no mature ghost-state discipline, no general heap model, and no interactive proof layer. \\
\bottomrule
\end{tabularx}
\end{center}

The strongest examples align with the language's first-order, loop-based style. These are not toy properties: they are exactly the obligations that make ordinary systems code fragile when written in an unchecked language. The implemented language is deliberately narrower than Dafny, Verus, or Liquid Haskell as a general verification environment. Its distinct contribution is that the properties it can express are carried into the compiled low-level representation and checked there independently.

\subsection{Comparison with verification-oriented languages}

The most appropriate comparison points are Dafny, Liquid Haskell, Verus, and ATS. Dafny is the closest imperative SMT-backed verifier, Liquid Haskell is the closest refinement-type automation baseline, Verus is the most relevant modern systems-language verification comparison, and ATS is the clearest historical predecessor for Veritas' broader goals.


\begin{table}[h]
  \centering
  \scriptsize
  \begin{tabular}{@{}p{0.43\textwidth} c c c c c@{}}
    \hline
    \textbf{Safety property / constraint class}         & \textbf{Veritas} & \textbf{Dafny} & \textbf{LH} & \textbf{Verus} & \textbf{ATS} \\
    \hline
    Array bounds safety                                 & \checkmark       & \checkmark     & partial                 & \checkmark     & \checkmark   \\
    Division-by-zero safety                             & \checkmark       & \checkmark     & \checkmark              & \checkmark     & \checkmark   \\
    Overflow / arithmetic side-condition reasoning      & \checkmark       & partial        & partial                 & \checkmark     & \checkmark   \\
    Preconditions and postconditions                    & \checkmark       & \checkmark     & \checkmark              & \checkmark     & \checkmark   \\
    Loop invariant reasoning                            & \checkmark       & \checkmark     & partial                 & \checkmark     & partial      \\
    Quantified specifications                           & \checkmark       & \checkmark     & partial                 & \checkmark     & partial      \\
    Sortedness / predicate-rich functional properties   & \checkmark       & \checkmark     & \checkmark              & \checkmark     & \checkmark   \\
    Ownership / aliasing discipline                     & \checkmark       & --             & --                      & \checkmark     & partial      \\
    Heap-allocated linked data structure invariants     & --               & \checkmark     & \checkmark              & partial        & partial      \\
    Recursive structural functional correctness proofs  & --               & \checkmark     & \checkmark              & \checkmark     & \checkmark   \\
    Concurrency / data-race freedom properties          & --               & --             & --                      & \checkmark     & --           \\
    Ghost-state-heavy protocol reasoning                & partial          & \checkmark     & partial                 & \checkmark     & \checkmark   \\
    Information-flow / noninterference-style properties & --               & partial        & partial                 & partial        & partial      \\
    Low-level preservation of checked constraints       & \checkmark       & --             & --                      & --             & --           \\
    Independent checking of compiled artefact           & \checkmark       & --             & --                      & --             & --           \\
    \hline
  \end{tabular}
  \caption{Comparison of the main safety properties and specification classes relevant to this dissertation. ``Partial'' indicates that the property is expressible, but less naturally or less directly in the style of programs considered here. LH abbreviates Liquid Haskell.}
  \label{tab:verification-safety-properties}
\end{table}


Table~\ref{tab:verification-safety-properties} makes the current limits of Veritas clearer: the implemented fragment is strong on array-heavy safety properties, arithmetic side conditions, contracts, and borrow-aware reasoning, but weak on heap-shaped invariants, recursive proof patterns, concurrency properties, and more protocol-oriented ghost-state verification. Where Veritas starts to diverge is in its ability to preserve the checked constraints in the fragments it does support below the source language.



The cross-system comparison suite contains six normalised tasks: preconditioned safe division, bounded array access with an offset, array initialisation with a loop invariant, binary search over a sorted collection, returning the head of a sorted array, and overflow-safe midpoint arithmetic. The suite is normalised by verified property: Dafny uses heap arrays, Veritas and Verus use fixed-size arrays naturally, and Liquid Haskell uses a more functional representation.

\begin{figure}[H]
  \centering
  \includegraphics[width=\linewidth]{figures/chapter4_annotation_burden.png}
  \caption{Annotation-to-code ratios across the normalised six-task
    cross-system comparison suite.}
  \label{fig:cross-system-annotation-burden}
\end{figure}

Within this small suite, Figure~\ref{fig:cross-system-annotation-burden} shows Veritas is competitive on annotation burden and is particularly light on the loop- and search-heavy tasks.



\section{Verification cost and performance}\label{sec:verification-perf}

This section asks how expensive those checks are, where the cost appears, and how Veritas should be compared with source-level verification systems that do not perform an explicit low-level re-checking stage.

\subsection{Veritas verification cost breakdown}

The first layer of cost evidence comes from the curated feature suite. These results separate source-to-DTAL compilation cost from the explicit DTAL verification stage.

\begin{center}
\small
\begin{tabularx}{0.88\textwidth}{>{\raggedright\arraybackslash}p{0.36\textwidth} r >{\raggedright\arraybackslash}X}
\toprule
\textbf{Metric} & \textbf{Mean} & \textbf{Worst case} \\
\midrule
Source-to-DTAL compilation & 27.110\,ms & \texttt{binary\_search}, 163.817\,ms \\
Explicit DTAL verification & 9.116\,ms & \texttt{binary\_search}, 55.209\,ms \\
Frontend SMT time & 13.097\,ms & Proof-heavy programs dominate \\
Verifier SMT time & 8.737\,ms & Proof-heavy programs dominate \\
\bottomrule
\end{tabularx}
\end{center}

For the current supported fragment, the mean explicit verifier cost is materially smaller than the mean frontend source-to-DTAL cost. The most expensive programs are also the proof-heavy ones: binary search mixes loop invariants, symbolic arithmetic, and sortedness-oriented reasoning, while simpler contract and borrow examples remain much cheaper.

The second layer of cost evidence comes from the verification-stress suite, a subset of the feature suite that acts as a proof-oriented stress suite of obligation-heavy programs.


\begin{table}[h]
  \centering
  \small
  \begin{tabular}{l r r r}
    \hline
    \textbf{Program}              & \textbf{Compile only (s)} & \textbf{Compile + verify (s)} & \textbf{Overhead ratio} \\
    \hline
    \texttt{array\_init}          & 0.0648                    & 0.2280                        & 3.521x                  \\
    \texttt{array\_assignment}    & 0.0890                    & 0.1089                        & 1.223x                  \\
    \texttt{quantifier\_showcase} & 0.1607                    & 0.1835                        & 1.142x                  \\
    \texttt{binary\_search}       & 0.2473                    & 0.2516                        & 1.017x                  \\
    \texttt{bubble\_sort}         & 0.0848                    & 0.1175                        & 1.385x                  \\
    \texttt{sortedness}           & 0.0283                    & 0.0394                        & 1.393x                  \\
    \hline
  \end{tabular}
  \caption{End-to-end wall-clock overhead of explicit DTAL verification on the
    verification-stress suite.}
  \label{tab:verification-stress-overheads}
\end{table}

Table~\ref{tab:verification-stress-overheads} shows that the cost of the independent verifier is highly program-dependent. \texttt{array\_init} is the clearest outlier, with a 3.521$\times$ overhead ratio, because it generates a particularly obligation-heavy array initialisation pattern.

\subsection{Cross-system verification-cost comparison}

The comparison suite supports a small empirical contrast with Dafny, Liquid Haskell, and Verus on six normalised tasks: safe division, bounded array access with an offset, array initialisation with a loop invariant, binary search, sorted-head extraction, and overflow-safe midpoint arithmetic. Repeated wall-clock timings were collected locally for Veritas in both \texttt{compile\_only} and \texttt{compile\_plus\_verify} modes, and for the other systems using their normal source-level verification commands.

\begin{figure}[H]
  \centering
  \includegraphics[width=0.86\linewidth]{figures/chapter4_comparison_suite_wallclock.png}
  \caption{Grouped wall-clock timings for the normalised six-task
    cross-system comparison suite. Bars show mean time and error bars show one
    standard deviation. The Veritas results were refreshed at 100 runs with 10
    warmups; the external-system results use the corresponding repeated local
    verification-command measurements.}
  \label{fig:comparison-suite-wallclock}
\end{figure}

Figure~\ref{fig:comparison-suite-wallclock} shows the resulting wall-clock times. Veritas averages around 0.038\,s in \texttt{compile\_only} mode and 0.047\,s in \texttt{compile\_plus\_verify} mode over 100 runs with 10 warmups. Dafny averages $\sim$1.016\,s, Liquid Haskell $\sim$1.227\,s, and Verus $\sim$0.488\,s, with \texttt{binary\_search} again the most expensive case for all systems. The comparison should not be read as a general speed ranking: Veritas verifies a deliberately smaller fragment, while the other systems run their full source-level verification pipelines. Its relevance is that the extra low-level checking stage remains small on the normalised tasks.

A key distinction is that Veritas incurs two separable verification costs, at the source and DTAL level, whereas most related systems primarily incur one.
The frontend-only view, approximated by \texttt{compile\_ms} and the frontend SMT metrics, is the fairest comparison point against source-level verification systems. The end-to-end view, measured by \texttt{compile\_plus\_verify} wall-clock timings and \texttt{verify\_ms}, is the relevant measure for the trust architecture.
The other systems report only source-level verification time because they do not attempt an analogous independently checkable low-level stage.

\section{Verifier functional correctness and reliability}\label{sec:verifier-correctness}

This section evaluates the verifier itself as an independent checking mechanism. The claim is not that the verifier has been proved correct in a formal metatheoretic sense, but that it behaves reliably in practice.

\begin{center}
\small
\begin{tabularx}{0.96\textwidth}{>{\raggedright\arraybackslash}p{0.22\textwidth} >{\raggedright\arraybackslash}p{0.20\textwidth} >{\raggedright\arraybackslash}X}
\toprule
\textbf{Evidence} & \textbf{Result} & \textbf{What it supports} \\
\midrule
\texttt{feature\_suite} & 28 / 28 accepted & Frontend compiler and independent physical-DTAL verifier agree on representative valid programs. \\
\texttt{negative\_suite} & 35 / 35 rejected & Source-level checking rejects representative unsafe programs across the main error modes (Table~\ref{tab:negative-suite-summary}). \\
\texttt{dtal\_tampering} & 14 / 14 rejected & The low-level verifier rejects representative corruptions introduced after source compilation. \\
\bottomrule
\end{tabularx}
\end{center}

The first two rows establish that the frontend and verifier agree on the supported source fragment: representative valid programs compile and verify, while representative unsafe programs are rejected before code generation. The third row is stronger evidence for Veritas' distinctive claim, because it checks what happens after a valid source program has already passed the frontend and an untrusted later phase corrupts the low-level artefact.

\subsection{Rejection of tampered DTAL}

The most distinctive reliability evidence is the DTAL tampering corpus, since it tests the trust boundary directly rather than only the source language. A valid source program is compiled to DTAL, the baseline DTAL is verified successfully, and then a small syntax-preserving mutation is applied to the low-level program. The verifier is expected to reject the tampered DTAL.

The current corpus contains 14 cases spanning return-signature mismatches, strengthened preconditions, contradictory assertions, positive and negative out-of-bounds accesses, arithmetic overflow, use-after-drop, ownership duplication, mutable-aliasing violations, and entry state or signature mismatches. The mutations are deliberately small, because this is the shape a backend bug can take.

\begin{center}
\small
\begin{tabularx}{0.96\linewidth}{@{}p{0.24\linewidth}X@{}}
\toprule
\textbf{Tamper} & \textbf{Why the verifier rejects it} \\
\midrule
\texttt{.returns int} changed to \texttt{.returns bool} & The return register's re-derived type no longer satisfies the declared function contract. \\
Safe array index changed to an out-of-bounds constant & The verifier can no longer prove the low-level load or store index is inside the array extent. \\
Duplicating move inserted after a live borrow & The ownership state now contains two usable aliases to a resource that should remain linear. \\
\bottomrule
\end{tabularx}
\end{center}

These mutations are rejected at exactly the point where an untrusted backend would otherwise be relied on.

This evidence shows that representative low-level corruptions are rejected even when the original source program was valid.


\section{Trust architecture and trusted computing base}\label{sec:tcb-eval}

This section returns to the dissertation's architectural claim: Veritas should not only verify useful programs, but do so with a smaller and more explicit trust boundary than a source-level-only verification pipeline. The main question here is what remains trusted once independent physical-DTAL verification is in place.

\subsection{Trusted computing base analysis}
In the current design, the frontend, lowering pipeline, register allocator, physical allocation pass, and most backend logic are intentionally untrusted: bugs in those components should be exposed as verification failure. Table~\ref{tab:tcb-summary} separates the measured verifier-side Rust core from the small trusted artefacts that remain below it on the Linux and bare-metal targets.

\begin{table}[ht]
  \centering
  \small
  \begin{tabular}{l l l c c}
    \hline
    \textbf{Trusted component} & \textbf{Form}         & \textbf{Size}             & \textbf{Linux} & \textbf{Unikernel} \\
    \hline
    DTAL verifier core         & Rust LOC              & 7372 / 34581 (21.3\%)     & \checkmark     & \checkmark         \\
    Direct encoder             & Rust LOC              & 328                       & \checkmark     & \checkmark         \\
    x86 instruction encoder    & Rust LOC              & 792                       & \checkmark     & \checkmark         \\
    Z3 solver                  & external tool         & large external dependency & \checkmark     & \checkmark         \\
    Runtime blob               & machine code          & 418 bytes                 & \checkmark     & --                 \\
    Multiboot bootstrap        & machine code          & 206 bytes                 & --             & \checkmark         \\
    Host OS / kernel           & execution environment & large external dependency & \checkmark     & --                 \\
    \hline
  \end{tabular}
  \caption{Decomposed trust architecture for the current Veritas
    implementation. The verifier-side Rust LOC measurement captures the main
    trusted implementation core; the remaining rows distinguish the additional
    trusted artefacts for Linux and bare-metal targets.}
  \label{tab:tcb-summary}
\end{table}

The verifier-side Rust core is \num{7372} lines out of \num{34581} non-test Rust lines, or 21.3\% of the counted implementation.
While this remains a substantial audit target in absolute terms, it does mean that most of the compiler remains outside the trusted path. The decomposition also shows how I/O changes the trust model. On Linux, adding I/O adds a small trusted runtime blob and its typed interface, while the verifier still re-checks all call sites independently. On the bare-metal target, that runtime layer is largely replaced by DTAL-verified Veritas code, leaving mainly the encoder and the 206-byte bootstrap beneath the verifier.

The important comparison is not only the number of trusted lines, but which lines they are. A conventional source-level verification pipeline still trusts later compilation and code generation to preserve the source-level result. A fully verified compiler instead tries to prove that the whole translation is semantics-preserving, but that proof is a large artefact in its own right. Veritas takes a more modest route: it trusts an independent checker for a lower-level language and a small encoder beneath it. That does not remove all trust, since Z3, the verifier implementation, and the final encoder still matter, but it gives a concrete and inspectable boundary around the code responsible for accepting executable artefacts.

\section{Threats to validity and limitations}

The results support the dissertation's main claims, but only for the current fragment and benchmark policy.

\begin{center}
\small
\begin{tabularx}{0.96\textwidth}{>{\raggedright\arraybackslash}p{0.28\textwidth} >{\raggedright\arraybackslash}X}
\toprule
\textbf{Threat} & \textbf{Scope of the claim} \\
\midrule
Benchmark coverage & Strongest evidence is for arrays, local mutation, contracts, loop invariants, quantifiers, borrowing, and structured numeric kernels. It is not evidence for heap-heavy data structures, arbitrary pointer manipulation, concurrency, or large module ecosystems. \\
Cross-system comparison & Dafny, Liquid Haskell, Verus, and Veritas differ in host language, proof style, and what counts as a natural encoding of a property. Veritas also pays both source-level checking cost and DTAL-level checking cost. \\
Verifier assurance & The feature suite, negative suite, and tampering corpus show coherent behaviour on the supported fragment, but do not prove verifier soundness or completeness. \\
\bottomrule
\end{tabularx}
\end{center}

\section{Summary against project objectives}

\componentblock{Source safety obligations are expressible}{Contracts, invariants, quantifiers, arrays, and borrowing examples are supported, within the current fragment. The remaining limitation is language coverage rather than the shape of the verification architecture.}
\componentblock{Safety evidence survives lowering}{TAST, TIR, DTAL annotations, and the binary-search trace show the relevant facts moving through the compiler rather than being erased. This is the key type-preservation evidence: the compiler emits proof-relevant structure that a later checker can inspect.}
\componentblock{Low-level output is checked independently}{\texttt{--verify-dtal}, physical-DTAL verification, and agreement on 28 / 28 valid programs show that the verifier can check compiled artefacts without using the source compiler. The trusted base still includes Z3, the verifier implementation, and the final encoder, but not the optimisation or lowering passes.}
\componentblock{Executable output is produced}{The toolchain emits Linux ELF binaries, uses a small runtime, supports a bare-metal target, and runs the PolyBench-style kernels. It is therefore an executable compiler prototype rather than only a formal model, although it is not a GCC-level optimiser.}
\componentblock{Compiler bugs can be exposed after compilation}{The DTAL tampering corpus rejects 14 / 14 representative low-level corruptions. This demonstrates the intended failure mode: compiler output that no longer justifies its safety annotations is rejected after compilation, before execution. It is empirical evidence rather than a machine-checked proof of verifier correctness.}

The evaluation therefore supports the dissertation's central claim: Veritas is more than a source-level refinement prototype. It is a working type-preserving compilation pipeline in which proof-relevant information is preserved to a low-level representation and independently re-checked there.

The core objectives are met: Veritas implements a DTAL-inspired trust architecture, verifies compiled low-level artefacts independently, supports a non-trivial source fragment, produces executable code, and gives a meaningful external practicality story through the translated PolyBench-style kernels. The partial objectives are scope-related rather than architectural: runtime performance and code quality are not yet competitive with mature systems compilers, and the supported fragment remains narrower than a mature systems language or general-purpose verification platform.

The strongest overall conclusion is architectural. Veritas occupies a distinct and well-evidenced point between source-level verification tools and foundational low-level proof systems. It combines automated SMT-backed source reasoning with an independently checkable DTAL layer, reducing the amount of compiler infrastructure that must be trusted for the core safety claim.
